\documentclass[11pt]{article}

% ---------------------------------------------------------------------------
% Sum-free + affine-cap achievement games (Nofil genus): paper scaffold.
% Structure + non-proof sections drafted from the repo notes (2026-07-07).
% The theorem-proof sections live in kernels/*.tex; each has first-pass proof
% text, with conic-localization explicitly framed as partial/open-frontier work.
%
% Every drafted paragraph block below is flagged
%   % DRAFT -- SUBJECT TO LEAD-AUTHOR REVISION
% and every specific numeric count carries a "% SOURCE:" comment.
% ---------------------------------------------------------------------------

\usepackage[utf8]{inputenc}
\usepackage[T1]{fontenc}
\usepackage{amsmath}
\usepackage{amssymb}
\usepackage{amsthm}
\usepackage{mathtools}
\usepackage{booktabs}
\usepackage[margin=1in]{geometry}
\usepackage[colorlinks=true,linkcolor=blue,citecolor=blue,urlcolor=blue]{hyperref}

% --- theorem environments (used by the kernel sections) ---
\theoremstyle{plain}
\newtheorem{theorem}{Theorem}
\newtheorem{lemma}{Lemma}
\newtheorem{corollary}{Corollary}
\newtheorem{proposition}{Proposition}
\newtheorem{conjecture}{Conjecture}
\theoremstyle{definition}
\newtheorem{definition}{Definition}
\newtheorem{example}{Example}
\theoremstyle{remark}
\newtheorem{remark}{Remark}

% --- common macros ---
\newcommand{\Znn}{\mathbb{Z}}
\newcommand{\Zn}[1]{\mathbb{Z}_{#1}}
\newcommand{\Fq}{\mathbb{F}_q}
\newcommand{\AG}{\mathrm{AG}}
\newcommand{\PG}{\mathrm{PG}}
\newcommand{\STS}{\mathrm{STS}}
\newcommand{\Grundy}{\mathcal{G}}       % Sprague-Grundy value

% ===========================================================================
\title{%
  % TODO(title): placeholder -- lead author to finalize.
  Achievement games in the Nofil genus:\\
  the sum-free game on $\Zn{n}$ and the cap game on $\AG(n,q)$}

\author{%
  % TODO(authors): placeholder author block.
  Tavis Rudd \and (co-authors TBD)}

\date{Draft \today}

\begin{document}
\maketitle

% ---------------------------------------------------------------------------
\begin{abstract}
% DRAFT -- SUBJECT TO LEAD-AUTHOR REVISION
We study impartial, normal-play achievement (building) games in the genus introduced by
Huggan, Huntemann, and Stevens under the name \emph{Nofil}~\cite{huggan2022nofil}: a move adds
a point to the built set as long as it does not complete a forbidden triple, and the last
player able to move wins. We determine the outcome for two structured families. First, for the
\emph{sum-free achievement game} on $\Zn{n}$ (build a sum-free subset of $\Zn{n}$; a move adds
$x$ keeping the set free of solutions of $a+b=c$), we prove that for every $n \geq 1$ the second
player wins if and only if $n \equiv 0, 1, 5 \pmod 6$. The associated sequence of Sprague--Grundy values
appears to be new. Second, for the \emph{cap achievement game} on $\AG(n,q)=\Fq^n$ (build a cap,
i.e.\ a set with no three collinear points), we prove that the second player wins for every $n$
and every prime power $q$; the $q=3$ case is the cap-set game, so this settles it in every
dimension---including $d=5$---with no computation. Translated into Nofil's language, the affine
theorem gives Nofil its first infinite family of Steiner triple systems with a determined
outcome (nim-value $0$ on $\STS(3^n)$ for all $n$), and the $\mathbb{F}_2$ sum-free data
supplies structured entries that cut against a $v \bmod 6$ nim-parity trend visible in the
computed Nofil data. We record extensive computational corroboration and a conic-localization
lemma directed at the open projective analogue.
\end{abstract}

% ===========================================================================
\section{Introduction}
\label{sec:intro}

% DRAFT -- SUBJECT TO LEAD-AUTHOR REVISION
An impartial \emph{achievement game} on a hypergraph is played by two players who alternately
choose previously unchosen vertices; a vertex is a legal move only while the set of chosen
vertices, together with the new vertex, remains free of any hyperedge. Under normal play the
last player able to move wins. Equivalently, the legal positions form the legal complex of an
impartial strong placement game, and residual positions are links in that
complex~\cite{faridi2019gamecomplexes,huntemann2019strongplacement}. Huggan, Huntemann, and
Stevens~\cite{huggan2022nofil} study this
game on the block hypergraph of a Steiner triple system under the name \emph{Nofil} (``Next One
to FIll is the Loser'': a move that would complete a block of already-chosen points is illegal).
They show, among other things, that deciding the outcome of Nofil positions on Steiner triple
systems is PSPACE-complete, and that once the available hypergraph collapses to a graph the game
is exactly Node-Kayles~\cite{sieben2023hypergraph}. General positions of this genus are therefore
intractable, which places \emph{theorems for structured families} at the tractable frontier.

% DRAFT -- SUBJECT TO LEAD-AUTHOR REVISION
This paper determines the outcome for two such families, both instances of the same genus.
\begin{itemize}
  \item \textbf{The sum-free game on $\Zn{n}$} (Section~\ref{sec:sumfree-zn}). A position is a
    sum-free subset $A \subseteq \Zn{n}$ (no $a+b=c$ with $a,b,c \in A$, and $a=b$ permitted so
    that $2a \in A$ is forbidden); a move adds an element keeping $A$ sum-free. This is Nofil's
    genus on the Schur $3$-uniform hypergraph of $\Zn{n}$, which is not a Steiner triple system.
    We prove a clean residue law: for every $n \geq 1$ the second player wins iff $n \equiv 0,1,5
    \pmod 6$.
  \item \textbf{The cap game on $\AG(n,q)$} (Section~\ref{sec:affine-cap}). A position is a cap
    in $\AG(n,q)=\Fq^n$ (no three collinear points); a move adds a point keeping it a cap. We
    prove that the second player wins for every $n$ and every prime power $q$. For $q=3$ this is
    the cap-set game and the result settles every dimension without computation.
\end{itemize}

% DRAFT -- SUBJECT TO LEAD-AUTHOR REVISION
We then draw out the consequences in Nofil's own terms (Section~\ref{sec:nofil-corollaries}).
Every $\AG(n,3)$ is a Steiner triple system $\STS(3^n)$, so the affine theorem gives Nofil its
first infinite family with a determined outcome; since these systems are vertex-transitive, their
nim-value lies in $\{0,1\}$ (Proposition~6 of~\cite{huggan2022nofil}), and a second-player win
pins it to exactly $0$. The $\mathbb{F}_2$ sum-free game coincides with the cap game on
$\PG(m,2)=\STS(2^{m+1}-1)$; its computed values give the projective column and a structured
observation about the $v \bmod 6$ nim-parity trend seen in the Nofil data. The projective cap
game in general remains open, and Section~\ref{sec:conic-localization} records a conic-localization
lemma directed at it.

% DRAFT -- SUBJECT TO LEAD-AUTHOR REVISION
\paragraph{Relation to prior work and novelty.} The game genus itself is published: our games
are Nofil~\cite{huggan2022nofil} on particular $3$-uniform hypergraphs, and we frame our
contributions as new theorems and infinite families \emph{inside} that genus rather than as a new
game. Within the genus, the sum-free achievement game on $\Zn{n}$ appears to be a previously
unstudied instance---a prior-art search surfaced no published game on sum-free sets, and the
resulting Grundy sequence has no match in the OEIS---while the cap game specializes the genus to
the two classical infinite Steiner-triple-system families (affine and projective). Two independent
cross-checks against the published Nofil data confirm the placement: the Fano plane
$\PG(2,2)=\STS(7)$ has nim-value $0$, matching our $\mathbb{F}_2$ result, and the unique
$\STS(9)=\AG(2,3)$ has nim-value $0$, matching the affine theorem's specialization.
% SOURCE: cross-checks and novelty verdict, notes/2026-07-07-nofil-connection.md and
% notes/2026-07-04-sumfree-game-theorem.md (revised 2026-07-07).

% ===========================================================================
\section{Preliminaries and conventions}
\label{sec:prelim}

% DRAFT -- SUBJECT TO LEAD-AUTHOR REVISION
We use standard impartial-game theory. All games here are finite, impartial, and played under
the \emph{normal-play} convention: the player facing a position with no legal move loses. For a
position $P$ we write $\Grundy(P)$ for its Sprague--Grundy (nim-) value; $P$ is a \emph{$P$-position}
(a second-player win) iff $\Grundy(P)=0$, and an \emph{$N$-position} (a first-player win) otherwise.
The Sprague--Grundy value of a disjunctive sum of independent components is the nim-sum (bitwise
XOR) of the components' values.

% DRAFT -- SUBJECT TO LEAD-AUTHOR REVISION
\begin{definition}[Achievement game on a hypergraph]
Let $H=(V,\mathcal{E})$ be a hypergraph. In the \emph{achievement game} on $H$, players alternately
add a vertex to a jointly built set $A \subseteq V$; adding $x \notin A$ is legal iff $A \cup \{x\}$
contains no hyperedge of $H$. Play starts from $A=\varnothing$; the last player to move wins.
Equivalently this is \emph{hypergraph Node-Kayles}: the reachable positions are the hyperedge-free
(``independent'') subsets of $H$, and terminal positions are the inclusion-maximal ones.
\end{definition}

% DRAFT -- SUBJECT TO LEAD-AUTHOR REVISION
\begin{remark}
For a $3$-uniform $H$ this is exactly Nofil~\cite{huggan2022nofil} on $H$ (a move completing a
hyperedge is illegal; normal play). When $H$ is the block hypergraph of a Steiner triple system,
it is Nofil in the original sense. Our two families use $H$ equal to the Schur triple hypergraph
of $\Zn{n}$ and the collinear-triple hypergraph of $\AG(n,q)$; the latter is a Steiner triple
system exactly when $q=3$.
\end{remark}

% DRAFT -- SUBJECT TO LEAD-AUTHOR REVISION
We record the objects and conventions used throughout.
\begin{itemize}
  \item \emph{Sum-free set.} $A \subseteq \Zn{n}$ is \emph{sum-free} if there is no solution of
    $a+b=c$ with $a,b,c \in A$; we allow $a=b$, so $2a=c$ is also forbidden. The element $0$ is
    never playable ($0+0=0$), so the game lives on $\Zn{n} \setminus \{0\}$.
  \item \emph{Cap.} $A \subseteq \AG(n,q)=\Fq^n$ is a \emph{cap} if it contains no three distinct
    collinear points. A line over $\Fq$ has $q$ points; for $q=3$, ``no three collinear'' is ``no
    full line,'' the classical cap-set condition, and $a,b,c$ collinear iff $a+b+c=0$.
  \item \emph{Symmetry.} Both games carry rich automorphism groups (affine maps; multiplications
    and negation on $\Zn{n}$). A hypergraph automorphism that is a fixed-point-free involution with
    no vertex adjacent to its image yields a mirror (pairing) strategy for the second player; this
    is the engine of both theorems and is developed in the respective sections.
  \item \emph{Vertex-transitivity.} If the automorphism group acts transitively on vertices, all
    opening moves are equivalent, so the root nim-value lies in $\{0,1\}$; this is Proposition~6
    of~\cite{huggan2022nofil} and applies to every $\AG(n,q)$ and every $\Zn{n}$.
  \item \emph{Decomposition.} The sum-free game decomposes as a disjunctive sum over the connected
    components of the residual ``armed'' Schur hypergraph, with $\Grundy$ the nim-sum of the
    components' values; we use this only as a solver lever and structural handle, not in the proofs.
\end{itemize}

% ===========================================================================
% --- The theorem-proof sections ---
% ===========================================================================

\section{The sum-free achievement game on \texorpdfstring{$\mathbb{Z}_n$}{Zn}}
\label{sec:sumfree-zn}

%TODO (lead author, F1a): statement + full proof of the mod-6 outcome theorem for the
%sum-free achievement game on Z_n -- Lemmas 1-4 (negation mirror; translation mirror by
%n/2; n/2 self-block; negation mirror with fixed extras) and the six-residue obstruction
%count. Source: notes/2026-07-04-sumfree-game-theorem.md. Do NOT draft proof text here.


\section{The cap achievement game on \texorpdfstring{$\mathrm{AG}(n,q)$}{AG(n,q)}}
\label{sec:affine-cap}

Let $\AG(n,q)$ denote the affine space $\Fq^n$, where $q$ is a prime power. A
\emph{cap} is a set of points containing no three distinct collinear points. The cap
achievement game starts from the empty cap and alternately adds points while preserving
the cap condition.

\begin{theorem}[Affine cap outcome law]
\label{thm:affine-cap}
For every $n\geq 1$ and every prime power $q$, the cap achievement game on $\AG(n,q)$
is a second-player win.
\end{theorem}

For $q=3$ this is the cap-set achievement game on $\mathbb F_3^n$, so the theorem settles
that game in every dimension without computation.

We use two elementary facts about affine mirrors. A \emph{cap automorphism} is a
permutation of $\Fq^n$ taking affine lines to affine lines; every invertible affine map is
one. Such a map preserves caps and legal moves.

\begin{lemma}[Line-parity lemma]
\label{lem:affine-line-parity}
Let $\sigma$ be a cap automorphism, let $A$ be a $\sigma$-symmetric cap, and let $\ell$
be a line with $\sigma(\ell)=\ell$. Suppose that on $\ell$, the map $\sigma$ acts
fixed-point-freely away from a set $F\subseteq \ell$, and that $A\cap F=\varnothing$.
If some $y\in \ell$ is a legal move at $A$, then $A\cap \ell=\varnothing$.
\end{lemma}

\begin{proof}
The set $A\cap\ell$ is $\sigma$-invariant and misses $F$, so it is partitioned into
2-element $\sigma$-orbits. Hence $|A\cap\ell|$ is even. Since $A\cup\{y\}$ is a cap and
$y\in\ell$, the line $\ell$ contains at most two points of $A\cup\{y\}$, so
$|A\cap\ell|\leq 1$. The only even possibility is $0$.
\end{proof}

\begin{lemma}[Mirror step]
\label{lem:affine-mirror-step}
Let $\sigma$ be an involutive cap automorphism, let $A$ be a $\sigma$-symmetric cap, and
let $y$ be a legal move at $A$ with $\sigma(y)\neq y$. Let $\ell$ be the line through
$y$ and $\sigma(y)$. If $A\cap\ell=\varnothing$, then $\sigma(y)$ is a legal reply and
$A\cup\{y,\sigma(y)\}$ is a cap.
\end{lemma}

\begin{proof}
If $\sigma(y)\in A$, then $y=\sigma(\sigma(y))\in \sigma(A)=A$, contradicting that $y$
is a move. Suppose $A\cup\{y,\sigma(y)\}$ contains three distinct collinear points. Any
such triple must contain $\sigma(y)$, since $A\cup\{y\}$ is already a cap. If the other
two points are in $A$, applying $\sigma$ gives a collinear triple in $A\cup\{y\}$, a
contradiction. If the triple is $\{y,\sigma(y),p\}$ with $p\in A$, then $p$ lies on
the line $\ell$, contradicting $A\cap\ell=\varnothing$.
\end{proof}

\begin{proof}[Proof of Theorem~\ref{thm:affine-cap}]
First suppose $q$ is even. If $q=2$, every subset is a cap and the board has $2^n$ points,
an even number, so the second player wins by any fixed pairing of the board.

Now let $q\geq 4$ be even. Choose a nonzero vector $v\in\Fq^n$ and define
\[
  \tau_v(x)=x+v.
\]
In characteristic $2$, $\tau_v$ is a fixed-point-free involution and an affine cap
automorphism. Let
\[
  \mathcal P=\{A\subseteq\Fq^n: A\text{ is a cap and }A=\tau_v(A)\}.
\]
The empty set lies in $\mathcal P$. If $A\in\mathcal P$ and $y$ is a legal move, let
$\ell=\{y+sv:s\in\Fq\}$ be the line through $y$ and $\tau_v(y)$. The translation
$\tau_v$ preserves $\ell$ and acts on it as $s\mapsto s+1$, fixed-point-freely. By
Lemma~\ref{lem:affine-line-parity}, $A\cap\ell=\varnothing$; by
Lemma~\ref{lem:affine-mirror-step}, $\tau_v(y)$ is a legal reply and the resulting cap
is again in $\mathcal P$. Lemma~\ref{lem:sumfree-responder} proves that the empty
position is $P$.

Now suppose $q$ is odd. The board has odd size, so a whole-board fixed-point-free mirror
cannot exist. Instead the second player builds a blocked fixed point. Let the first
player open at $a$. The second player chooses any $b\neq a$, which is legal because two
points always form a cap. Let
\[
  c=\frac{a+b}{2}.
\]
Then $a,b,c$ are distinct collinear points, so once $a$ and $b$ have been played, $c$ is
not legal for the rest of the game. Let
\[
  \sigma_c(x)=2c-x
\]
be the point reflection through $c$. This is an affine involution, fixes exactly $c$, and
swaps $a$ and $b$. Define
\[
  \mathcal P=\{A\subseteq\Fq^n: A\text{ is a cap},\ A=\sigma_c(A),\ \{a,b\}\subseteq A,
    \ c\notin A\}.
\]
The position $\{a,b\}$ lies in $\mathcal P$. If $A\in\mathcal P$ and $y$ is legal, then
$y\neq c$, since $c$ is blocked by $a,b$. Thus $\sigma_c(y)\neq y$. The line $\ell$
through $y$ and $\sigma_c(y)$ passes through $c$ and is $\sigma_c$-invariant; on
$\ell\setminus\{c\}$, the reflection is fixed-point-free. Since $A\cap\{c\}=\varnothing$,
Lemma~\ref{lem:affine-line-parity} gives $A\cap\ell=\varnothing$, and
Lemma~\ref{lem:affine-mirror-step} makes $\sigma_c(y)$ a legal reply. The new position
is again in $\mathcal P$. Therefore $\{a,b\}$ is a $P$-position by
Lemma~\ref{lem:sumfree-responder}. Since the second player can reach such a position
after every opening $a$, the empty position is $P$.
\end{proof}

The two cases explain why affine space is easier than the projective analogue. In even
characteristic, translations give a fixed-point-free pairing of the whole board. In odd
characteristic, the unavoidable fixed point of a reflection is placed on the opening secant,
where it is permanently blocked. Projective spaces have neither translations nor such a
one-point fixed set, which is why the projective cap game remains open.


\section{Corollaries for Nofil on Steiner triple systems}
\label{sec:nofil-corollaries}

%TODO (lead author, F1c): the Nofil corollaries -- (i) AG(n,3) as an STS gives the first
%infinite determined Nofil family, nim-value 0 on STS(3^n) for all n (via vertex-transitivity,
%Prop. 6 of Huggan-Huntemann-Stevens); (ii) the PG(m,2) = STS(2^{m+1}-1) sum-free column and
%the v mod 6 nim-parity observation. Prove the convention equivalence carefully (the one place
%a silent error kills the paper). Source: notes/2026-07-07-nofil-connection.md. Do NOT draft here.


\section{Projective cap games and conic localization}
\label{sec:conic-localization}

We close with a partial result for the projective analogue. This section is not a proof that
the projective cap game is always a second-player win; it isolates the finite-geometric
structure of the local obstruction that remains after the affine theorem.

\paragraph{The residual grid game.}
After two projective opening points are placed on the line at infinity, the remaining affine
points may be represented as a grid $\Fq\times\Fq$. A residual position is a set of cells
that is both a partial permutation (no two cells in a common row or column) and an affine cap
(no three cells collinear in $\AG(2,q)$). A move adds a cell preserving both conditions.
The row and column constraints encode the two already-played directions at infinity.

Embed $(r,c)\in\Fq^2$ as $(r:c:1)\in\PG(2,q)$ and let
\[
  a=(1:0:0),\qquad b=(0:1:0)
\]
be the two direction points. Lines through $a$ are affine columns, lines through $b$ are
affine rows, and the line $ab$ is the line at infinity.

\begin{lemma}[The associated 5-arc]
\label{lem:grid-five-arc}
Let $S_3=\{(r_i,c_i):1\leq i\leq 3\}$ be a size-$3$ residual grid position. Then
$S_3\cup\{a,b\}$ is a $5$-arc in $\PG(2,q)$.
\end{lemma}

\begin{proof}
A collinear triple containing both $a$ and $b$ would lie on the line at infinity, but all
cells of $S_3$ have affine coordinate $z=1$. A triple containing $a$ and two cells would
put those two cells in one column; a triple containing $b$ and two cells would put them in
one row. Both are forbidden by the partial-permutation condition. A triple consisting of
the three cells is forbidden by the affine-cap condition.
\end{proof}

\begin{theorem}[Conic localization]
\label{thm:conic-localization}
Let $S_3$ be a size-$3$ residual grid position. Then there is a unique nondegenerate conic
$\mathcal C\subseteq\PG(2,q)$ through $S_3\cup\{a,b\}$. Its affine part has the form
\[
  \mathcal C_{\rm aff}
  =\{(r,c)\in\Fq^2:(r-\rho)(c-\alpha)=B\}
\]
for uniquely determined $\rho,\alpha\in\Fq$ and $B\in\Fq^\ast$. In particular,
$\mathcal C_{\rm aff}$ has exactly $q-1$ cells, one in each row $r\neq\rho$ and one in
each column $c\neq\alpha$.

Moreover, every cell of $\mathcal C_{\rm aff}\setminus S_3$ is a legal extension of
$S_3$. The set $\mathcal C_{\rm aff}$ is itself a residual grid position; if $q$ is odd,
it is inclusion-maximal and has even size $q-1$.
\end{theorem}

\begin{proof}
Write a conic as the zero set of a nonzero quadratic form
\[
Q(r,c,z)=\lambda r^2+\mu c^2+\gamma z^2+\delta rc+\epsilon rz+\zeta cz.
\]
The conditions $Q(a)=Q(b)=0$ force $\lambda=\mu=0$. If $\delta=0$, then
$Q=z(\epsilon r+\zeta c+\gamma z)$, so the three affine cells of $S_3$ would lie on one
line, contradicting the cap condition. Thus $\delta\neq 0$; scale so that $\delta=1$.
The affine equation is then
\[
  rc+\epsilon r+\zeta c+\gamma=0.
\]
The three cell conditions are a linear system in $(\epsilon,\zeta,\gamma)$ with coefficient
rows $(r_i,c_i,1)$. Its determinant is the collinearity determinant of the three cells,
nonzero because $S_3$ is an affine cap. Hence the conic is unique. Completing the product
gives
\[
  rc+\epsilon r+\zeta c+\gamma=(r-\rho)(c-\alpha)-B
\]
with $\rho=-\zeta$, $\alpha=-\epsilon$, and $B=\rho\alpha-\gamma$.

We have $B\neq 0$. If $B=0$, the affine equation factors as
$(r-\rho)(c-\alpha)=0$, the union of one row and one column; three cells lying in this
union force two to share a row or a column. The same factorization argument also shows
nondegeneracy: if the projective conic factored into two lines, the five arc points would
lie on two lines, hence three of them on one line, contradicting
Lemma~\ref{lem:grid-five-arc}.

Since $B\neq 0$, each row $r\neq\rho$ contains the unique conic cell
\[
  c=\alpha+\frac{B}{r-\rho},
\]
and no cell occurs in row $\rho$; similarly, each column $c\neq\alpha$ occurs exactly once.
Thus $\mathcal C_{\rm aff}$ has $q-1$ cells and is a partial permutation.

Let $x\in\mathcal C_{\rm aff}\setminus S_3$. The row and column uniqueness just proved show
that $x$ uses none of the three rows or columns already used by $S_3$. If $x$ and two cells
of $S_3$ were collinear, a line would meet the nondegenerate conic $\mathcal C$ in three
distinct points, impossible. Hence $S_3\cup\{x\}$ is a legal residual position. The same
two-point line bound shows that $\mathcal C_{\rm aff}$ itself is an affine cap, so it is a
residual grid position.

Assume now that $q$ is odd. Let $x=(r,c)\notin\mathcal C_{\rm aff}$. If $r\neq\rho$, then
row $r$ already contains a conic cell, so $x$ is illegal. If $r=\rho$ and $c\neq\alpha$,
then column $c$ already contains a conic cell, so $x$ is illegal. The only remaining cell is
the center $(\rho,\alpha)$. In translated coordinates $t=r-\rho$, $s=c-\alpha$, the conic
cells are
\[
  P_t=(t,B/t),\qquad t\in\Fq^\ast.
\]
For any $t\in\Fq^\ast$, the two points $P_t$ and $P_{-t}$ are distinct, and the center is
on the line through them. Adding the center therefore creates a collinear triple. Thus
every off-conic cell is illegal, and $\mathcal C_{\rm aff}$ is maximal.
\end{proof}

\begin{remark}[Even characteristic]
\label{rem:even-conic-nucleus}
The maximality conclusion is genuinely odd-characteristic. If $q$ is even, the center
$(\rho,\alpha)$ is the nucleus of the conic: no secant through two distinct conic cells
passes through it, and $\mathcal C_{\rm aff}\cup\{(\rho,\alpha)\}$ is a larger residual
position. This is harmless for the main theorem, since even-characteristic affine games are
already settled by the translation mirror of Section~\ref{sec:affine-cap}.
\end{remark}

\begin{corollary}[On-conic extension count]
\label{cor:on-conic-count}
Every size-$3$ residual grid position has exactly $q-4$ legal extensions on its conic. If
one also uses the total extension count $q^2-9q+21$, the number of off-conic legal extensions
is $(q-5)^2$.
\end{corollary}

\begin{proof}
The conic has $q-1$ affine cells, three of which are the cells of $S_3$, and
Theorem~\ref{thm:conic-localization} says that all remaining conic cells are legal. The
off-conic formula is the subtraction
\[
  q^2-9q+21-(q-4)=q^2-10q+25=(q-5)^2.
\]
\end{proof}

\begin{lemma}[Multiplicative conic reflections]
\label{lem:psi-u}
Work in translated coordinates $t=r-\rho$, $s=c-\alpha$, so that
$\mathcal C_{\rm aff}=\{(t,B/t):t\in\Fq^\ast\}$. For $u\in\Fq^\ast$ define
\[
  \psi_u(t,s)=\left(\frac{u}{B}s,\frac{B}{u}t\right).
\]
Then $\psi_u$ is an involutive grid symmetry, it maps $\mathcal C_{\rm aff}$ to itself, and
on the conic parameter it acts by $t\mapsto u/t$. Its fixed-point set is the line
$s=(B/u)t$; for $q$ odd this line meets the conic in two cells if $u$ is a square and in no
cells if $u$ is a nonsquare.
\end{lemma}

\begin{proof}
The displayed map is linear, invertible, and squares to the identity. It exchanges the row
and column families, so it preserves the partial-permutation condition, and being linear it
preserves affine collinearity. On a conic point,
\[
  \psi_u(t,B/t)=\left(u/t,\,B/(u/t)\right),
\]
so the parameter changes by $t\mapsto u/t$. The fixed equations are equivalent to
$s=(B/u)t$. Intersecting this line with the conic gives $t=u/t$, i.e. $t^2=u$, which has
two or zero solutions over $\Fq$ when $q$ is odd according as $u$ is or is not a square.
\end{proof}

\begin{conjecture}[On-conic escape]
\label{conj:on-conic-escape}
For odd $q$, every size-$3$ residual grid position has at least one legal extension on its
conic that is a $P$-position of the residual grid game.
\end{conjecture}

Conjecture~\ref{conj:on-conic-escape} implies the local escape statement used in the frame
reduction for $\PG(2,q)$, and hence would imply the projective-plane second-player win in
that reduction. It is strictly stronger than the escape statement: a class with all on-conic
extensions $N$ but an off-conic $P$ extension would refute the conjecture without refuting the
projective outcome.

The evidence is exhaustive through the current wall: $q=5,7,9,11,13,17,19$, including the
non-prime field $q=9$. The data also rule out several simpler laws. In particular, not every
on-conic extension is a $P$-position, the value is not determined by a fixed quadratic-character
formula in the conic parameters, and parity of the number of maximal completions does not decide
the on-conic value. The useful reduction is nevertheless substantial: the witness search drops
from the full two-dimensional grid to the one-dimensional conic parameter set
$\Fq^\ast\setminus S_3$ modulo the rotations $t\mapsto kt$ and the reflections of
Lemma~\ref{lem:psi-u}.

For deeper steering, the natural finite-geometry language is the conic-stabilizer action of
$PGL(2,q)$. Hollmann and Xiang describe the associated coherent configuration and association
schemes in terms of cross-ratio relations~\cite{hollmann2005pglconic}. Their intersection numbers
may provide a way to count response cells in a prescribed orbital relation, potentially improving
on coarse row/column reservoir estimates when one needs an exact positive count.


% ===========================================================================
\section{Computational results}
\label{sec:computational}

% DRAFT -- SUBJECT TO LEAD-AUTHOR REVISION
The theorems above are proved uniformly and no machine step is load-bearing; the computations in
this section are corroboration, and they also produced the new integer sequence.

% DRAFT -- SUBJECT TO LEAD-AUTHOR REVISION
\paragraph{The sum-free sequence on $\Zn{n}$.} A multiplier-quotient Grundy solver (Rust) computed
$\Grundy(\varnothing)$ for $\Zn{n}$; the values were validated against a direct brute-force solver
for small $n$, and the outcome law was confirmed with zero exceptions.
The first terms $a(n)=\Grundy(\varnothing)$, $n=1,\dots,27$, are
% SOURCE: notes/2026-07-04-sumfree-oeis-draft.md (b-file n=1..65) and
% notes/2026-07-04-sumfree-capset-game.md (Result 1 table).
\[
  0,\,1,\,1,\,2,\,0,\,0,\,0,\,2,\,1,\,1,\,0,\,0,\,0,\,2,\,2,\,3,\,0,\,0,\,0,\,2,\,1,\,3,\,0,\,0,\,0,\,2,\,1,\dots
\]
The sequence was computed through $n=65$; all values through $n=65$ lie in $\{0,1,2,3\}$, and the
sequence is not eventually periodic---only the \emph{outcome} (whether $a(n)=0$) is periodic mod
$6$. A direct OEIS lookup of the initial terms returns no match, so the sequence appears to be new.
% SOURCE: computed through n=65 and OEIS no-match, notes/2026-07-04-sumfree-oeis-draft.md;
% values validated vs brute-force for small n, notes/2026-07-04-sumfree-capset-game.md.

% DRAFT -- SUBJECT TO LEAD-AUTHOR REVISION
Two further checks support the proof mechanism. The mirror steps were verified to have zero
counterexamples over all reachable symmetric positions for $n \leq 36$, and the outcome law was
confirmed for every $n$ from $5$ to $65$ without exception.
% SOURCE: 0 mirror-breakers n<=36 and law confirmed to n=65, notes/2026-07-04-sumfree-game-theorem.md.
Independently, the disjunctive-sum decomposition (nim-sum over connected components of the residual
armed Schur hypergraph) was checked with zero mismatches over roughly $70{,}000$ reachable positions
for $n \in \{12,15,17,18,20,22,24,30\}$.
% SOURCE: ~70,000 positions, listed n, 0 mismatches, notes/2026-07-04-sumfree-capset-game.md.

% DRAFT -- SUBJECT TO LEAD-AUTHOR REVISION
\paragraph{The cap game on $\AG(n,q)$.} Brute-force Sprague--Grundy solves confirm the second-player
win ($\Grundy=0$) in every case attempted: $\AG(2,3)$ and $\AG(3,3)$ ($q=3$); $\AG(1,5),\AG(2,5)$ and
$\AG(1,7),\AG(2,7)$ (odd prime $q$); $\AG(2,9)$ (odd non-prime $q$, $81$ points, $410{,}575$ reachable
states); and $\AG(1,4),\AG(2,4)$ and $\AG(1,8),\AG(2,8)$ (characteristic $2$).
% SOURCE: enumerated cases incl. AG(2,9) 81 pts / 410,575 states, notes/2026-07-04-capset-game-theorem.md.
For $q=3$ the four-dimensional case was solved via a full $\mathrm{AGL}(4,3)$ canonical quotient: the
outcome-only solve returns a $P$-position over $8191$ classes, and the full Grundy solve returns
$\Grundy(\AG(4,3))=0$ over $9908$ inequivalent caps (of all sizes $0$--$20$), the two agreeing and
reproduced by two independent binaries.
% SOURCE: AGL(4,3) quotient, 8191 outcome classes / 9908 Grundy classes, notes/2026-07-04-sumfree-capset-game.md.
The explicit mirror strategies were validated to beat every first-player line of play: the reflection
mirror for $\AG(2,3),\AG(2,5),\AG(2,7),\AG(3,3)$ and the translation mirror for $\AG(2,4),\AG(2,8)$.
% SOURCE: strategy validation runs, notes/2026-07-04-capset-game-theorem.md.

% DRAFT -- SUBJECT TO LEAD-AUTHOR REVISION
\paragraph{Cross-checks against the Nofil data.} The two smallest classical Steiner triple systems
match: $\STS(7)$ (the Fano plane $=\PG(2,2)$) has nim-value $0$, agreeing with the $\mathbb{F}_2$
sum-free result, and the unique $\STS(9)=\AG(2,3)$ has nim-value $0$, agreeing with the affine
theorem---both as computed in~\cite{huggan2022nofil}. For the projective column, the
$\mathbb{F}_2$ sum-free game equals the cap game on $\PG(m,2)=\STS(2^{m+1}-1)$ and is covered by
the characteristic-$2$ projective theorem. By contrast the generic Nofil data are irregular---for
example $\STS(13)$ has
nim-value $1$, most $\STS(15)$ and $\STS(21)$ have nim-value $0$ with values up to $3$ occurring,
and $\STS(25)$ shows mixed values---so the uniform $P$ on our infinite families is special rather
than generic.
% SOURCE: STS(7),STS(9) matches; PG(m,2) characteristic-2 projective theorem; generic STS calibration
% (STS 13/15/21/25) -- notes/2026-07-07-nofil-connection.md (generic figures from Huggan-Huntemann-Stevens).

% ===========================================================================
\section{Related work}
\label{sec:related}

% DRAFT -- SUBJECT TO LEAD-AUTHOR REVISION
\paragraph{The Nofil genus.} Huggan, Huntemann, and Stevens~\cite{huggan2022nofil} introduced Nofil
on Steiner triple systems, computed nim-values for all systems of small order, proved
PSPACE-completeness of the outcome problem (via a Node-Kayles reduction), and observed the
Node-Kayles endgame. Their program is exhaustive to small order plus samples at larger orders, and
no infinite family had a determined outcome; a $2025$ follow-up~\cite{huntemann2025minimal} develops
the associated point-deletion / graph-embedding machinery. Our results live in this genus.

% DRAFT -- SUBJECT TO LEAD-AUTHOR REVISION
\paragraph{Impartial games on hypergraphs and groups.} Sieben's impartial hypergraph
games~\cite{sieben2023hypergraph} give the general achievement/building-game framework; our games
are the achievement game on specific $3$-uniform hypergraphs. The same positions are legal
complexes of impartial strong placement games, so link-of-a-face language gives the clean
description of residual games~\cite{faridi2019gamecomplexes,huntemann2019strongplacement}. A
closely related active line is the
impartial \emph{group generation} achievement games of Benesh, Ernst, and
Sieben~\cite{benesh2016symmetric,benesh2018nilpotent}, where the winning condition is generating the
group rather than staying sum-free or cap---a framework relative with a different objective.

% DRAFT -- SUBJECT TO LEAD-AUTHOR REVISION
\paragraph{Games on \textsc{Set}.} Several games share the $\AG(n,3)$ line structure but are distinct
from ours. \emph{Anti-Set}~\cite{clark2016antiset} is a mis\`ere game with separate hands and an
explicit strategy (no nim-values). The impartial \textsc{Set} game of Uiterwijk and
Hufkens~\cite{uiterwijk2023impartialset} is a \emph{removal} game (players take lines out of the
layout), whereas our players \emph{build} a cap by adding points---a different move set on the same
line hypergraph. Lampis and Mitsou~\cite{lampis2014setcomplexity} analyze the computational
complexity of the \textsc{Set} family. None of these compute the achievement-game nim-values studied
here.

% DRAFT -- SUBJECT TO LEAD-AUTHOR REVISION
\paragraph{Extremal cap sets and sum-free sets.} The cap game is adjacent to, but distinct from, the
extremal cap-set problem (the maximum size of a cap): Croot, Lev, and
Pach~\cite{croot2017progression} and Ellenberg and Gijswijt~\cite{ellenberg2017capset} bound the
maximum size for $q=3$, and program search has improved low-dimensional lower
bounds~\cite{romeraparedes2024funsearch}. These concern maximum size; our game concerns the outcome,
and the theorem says nothing about cap sizes (the game's terminals are inclusion-\emph{maximal}, not
\emph{maximum}, caps). Likewise the structure of sum-free and maximal-sum-free sets---classical
material of Wallis, Street, and Wallis~\cite{wallis1972combinatorics}, Diananda and
Yap~\cite{diananda1969maximal}, and Green and Ruzsa~\cite{green2005sumfree}---describes the terminal
positions of the $\Zn{n}$ game but not its outcome. The umbrella framework is Node-Kayles, whose
nim-value sequences have been tabulated for several graph families~\cite{wong2024nimber}.

% ===========================================================================
\section{Open problems}
\label{sec:open}

% DRAFT -- SUBJECT TO LEAD-AUTHOR REVISION
\begin{enumerate}
  \item \textbf{The projective cap game.} Is the cap game on $\PG(m,q)$ always a second-player win?
    It is computed to be $P$ in every small case, but the affine proof does not transfer:
    $\PG(m,q)$ has no translations and the board size $(q^{m+1}-1)/(q-1)$ has no uniform parity, so
    neither mirror drops out. Section~\ref{sec:conic-localization} is a partial step; a general
    certificate is open.
  \item \textbf{Even-dimensional odd projective spaces.} The binary projective column is no longer
    the open edge: $\PG(m,2)$ is covered by the characteristic-$2$ projective theorem.  The
    remaining higher-dimensional projective gap starts with even projective dimensions over odd
    fields, such as $\PG(4,3)$ and its analogues.
    % SOURCE: corrected projective-gap status, notes/2026-07-09-stepping-stone-deliverables-proposal.md.
  \item \textbf{The abelian generalization of the sum-free law.} The mod-$6$ law is the cyclic case
    of a criterion for the sum-free game on any finite abelian group, proved when the Sylow-$3$
    subgroup is cyclic or the $2$-rank is at least $2$; the slice with $2$-rank $\leq 1$ and $3$-rank
    $\geq 2$ (never occurring in $\Zn{n}$) remains open.
    % SOURCE: abelian criterion status, notes/2026-07-04-sumfree-game-theorem.md (generalization section).
  \item \textbf{Structure of the $N$-side nim-values.} The sum-free outcome is periodic mod $6$, but
    the nim-values themselves are irregular on the first-player-win residues; no closed form is known.
  \item \textbf{Higher-degree variants.} ``No $t$ collinear'' for $t>3$, or higher-flat achievement
    games, are not covered by the mirror argument (a mirror can create a longer forbidden flat) and
    are open.
  \item \textbf{Complexity of structured positions.} Nofil is PSPACE-complete on Steiner triple
    systems in general~\cite{huggan2022nofil}; the complexity of deciding positions of the
    $\AG(n,q)$ and $\PG(2,q)$ games specifically is open.
\end{enumerate}

% ===========================================================================
\bibliographystyle{plain}
\bibliography{refs}

\end{document}
